\section{平均値}\label{節:関数の平均}

\sectionref{節:定積分性質}で導入した関数の平均値を,温度のように時間とともに連続的に変わる量に使ってみよう.その前に,平均を積分で表す理由を振り返っておく.

$n$ 個の数 $u_1,\ldots,u_n$ の平均は,和を個数で割ったものである.関数 $f$ の区間 $[a,b]$ における平均も,区間から等間隔に $n$ 個の値を取り,その平均の極限として考える.\,$\Delta x=\frac{b-a}{n}$ と置けば,\,$\frac1n=\frac{\Delta x}{b-a}$ であるから,
\[
\frac1n\sum_{k=1}^n f(a+k\Delta x)
 =\frac{1}{b-a}\sum_{k=1}^n f(a+k\Delta x)\Delta x
\]
となる.右辺の和はリーマン和なので,\,$f$ が連続なら,\,$n\to\infty$ として
\begin{equation}\label{式:関数の平均値}
\overline f=\frac{1}{b-a}\int_a^b f(x)\dd x
\end{equation}
を得る.\sectionref{節:定積分性質}ではこれを $f_{\mathrm{av}}$ と書いたが,本章では $\overline f$ と書く.\,$f\geq0$ の場合,底辺 $b-a$,高さ $\overline f$ の長方形は,曲線の下の部分と同じ面積になる.\,$f$ が負の値をとっても,平均は同じ式で定める.

\begin{example}{変化する温度の平均}
時刻 $0\leq t\leq12$ の温度を $T(t)=20+6\sin\frac{\pi t}{12}$ とする.この時間帯の平均温度を求めよ.
\end{example}

\paragraphbf{方針}
式\eqref{式:関数の平均値}に従って,\,$T(t)$ を $0$ から $12$ まで積分し,時間の長さ $12$ で割る.

\paragraphbf{解答}
$20+\frac{12}{\pi}$

\paragraphbf{解法}
\begin{align*}
\overline T
 &=\frac1{12}\int_0^{12}\left(20+6\sin\frac{\pi t}{12}\right)\dd t\\
 &=\frac1{12}\left[20t-\frac{72}{\pi}\cos\frac{\pi t}{12}\right]_0^{12}\\
 &=\frac1{12}\left\{240-\frac{72}{\pi}(-1-1)\right\}\\
 &=20+\frac{12}{\pi}.
\end{align*}
始めと終わりの温度はともに $20$ だが,途中の温度は $20$ より高い.両端の温度の平均 $20$ は,時間全体の平均にはならないのである.

\paragraphbf{平均値をとる点}
\sectionref{節:定積分性質}で見たように,連続な $f$ の最小値を $m$,最大値を $M$ とすると $m\leq\overline f\leq M$ であり,ある $c\in[a,b]$ で $f(c)=\overline f$ となる(積分の平均値の定理).ただし,その $c$ が区間の中央にあるとは限らず,1つとも限らない.上の例題では $\sin\frac{\pi t}{12}=\frac2\pi$ を満たす $t$ が $0<t<12$ に2つあり,どちらも中央の $t=6$ とは異なる.

\paragraphbf{重み付き平均}
各点を均等に数えない平均もある.非負の連続関数 $w$ が $\int_a^b w(x)\dd x>0$ を満たすとき,\,$w$ を重みとする\textbf{重み付き平均}を
\[
\frac{\int_a^b f(x)w(x)\dd x}{\int_a^b w(x)\dd x}
\]
で定める.分母は重みの総量であり,\,$w=1$ とすれば式\eqref{式:関数の平均値}に戻る.次節で扱う重心は,位置の座標を面積や長さを重みとして平均したものである.
